\subsection{The graded loop equations: three spaces that must be distinguished}

The manuscript anchors are \S41, folio 315 (cote 140-3, scan 54),
\S42, folio 331 (scan 85), and its explicit-solution question at
folio 349 (scan 120).  The computations below answer the added graded
questions while keeping their profinite lifting problem explicit.

We fix a prime $\ell$, put $k=\mathbb Q_\ell$, and use positive word weight,
so that the corresponding motivic weight is $-2n$.  Write
$F^{(\ell)}=\langle X,Y\rangle_{\mathrm{pro}\text{-}\ell}$,
$Z=(XY)^{-1}$, and $\mathfrak L=\operatorname{Lie}_k\langle x,y\rangle$,
with $x=\log X$, $y=\log Y$ in unipotent completion.  In the associated
graded Lie algebra, and only there, $z=-x-y$.  In the filtered Lie algebra
$\log Z=-\operatorname{BCH}(x,y)$.
This distinction changes the answer to the dimension question.

\begin{claim}{B.g1}{Proved}{The precise leading-term map and its limitation}
Restrict normalized loop automorphisms to $\lambda=1$, writing
$\alpha_f(X)=X$, $\alpha_f(Y)=f^{-1}Yf$ with $f$ in the closed commutator
subgroup.  Give this group the filtration in which $f$ has lower-central
weight at least $n$.  For $n\geq2$ its degree-$n$ leading term is
\[
 \psi=(\log f)_n\in\mathfrak L_n,\qquad
 D_\psi(x)=0,\quad D_\psi(y)=[y,\psi].
\]
This gives an injective map on the associated graded of the normalized
pro-$\ell$ group, after tensoring with $k$, into outer special derivations.
If the outer automorphism commutes with the geometric $\mathfrak S_3$,
its image is $\mathfrak S_3$-invariant.  Its bracket is the Ihara bracket
\[
 \{\psi,\phi\}=D_\psi(\phi)-D_\phi(\psi)+[\psi,\phi].
\]
We do not identify the image of this map with the entire space of
invariant special derivations.
\end{claim}
% anchor:B.g1
\manuscriptanchor{\S 42; folios 331, 349; cote 140-3 scans 85, 120}

\begin{proof}
The graded Lie algebra of a free pro-$\ell$ group for the closed lower
central series is the free Lie algebra over $\mathbb Z_\ell$; its
rationalization is $\mathfrak L$.  One may obtain this description
from the integral Lyndon basis and the Magnus embedding.  Its
compatibility with the rational unipotent completion is supplied by
\citet[Theorem~A.6 and Corollary~A.10]{HM2003}.
In logarithmic coordinates,
$\alpha_f(y)=\exp(-\operatorname{ad}\log f)y
=y+[y,\psi]+\text{terms of degree greater than }n+1$.
A Lie element of degree greater than one commuting with $y$ is zero:
its associative expansion would be a polynomial in $y$, whereas the
ordinary free Lie algebra on one generator has no such degree.
Thus the normalized leading term is recovered from the automorphism.
Changing a general lift of an outer automorphism adds an inner
derivation.  The normalized lift has no ambiguity in degrees $n>1$,
since a possible inner change must centralize $x$ and hence vanish in
that degree.

Preservation of the conjugacy class of $Z$ implies that the leading
variation of $z$ is $[c,z]$ for some $c\in\mathfrak L_n$.
The geometric permutations of $X,Y,Z$ act on the associated graded by
permuting $x,y,z$; all higher terms in $\log Z$ disappear in this
leading comparison.  Commutation modulo inner automorphisms therefore
implies invariance of the outer derivation class.  Finally, applying
$[D_\psi,D_\phi]$ to $y$ and using Jacobi gives
$[D_\psi,D_\phi](y)=[y,\{\psi,\phi\}]$.
These arguments establish the inclusion and the bracket; they contain
no assertion that an arbitrary graded derivation has a compatible
integral lift at every subsequent degree.
\end{proof}

For clarity we define the exact formal system we compute.  Let
$\mathcal H_n\subset\mathfrak L_n$ be the solution space of
\begin{align}
 \psi(x,y)+\psi(y,x)&=0,\label{eq:graded-two}\\
 \psi(x,y)+\psi(y,z)+\psi(z,x)&=0,\qquad z=-x-y.
 \label{eq:graded-three}
\end{align}
The normalization $f\in(F^{(\ell)})'$ removes $\mathcal H_1$.
Let $W^{\mathrm{lac}}_n$ denote the verified graded loop-and-symmetry
system: the $\mathfrak S_3$-invariant outer derivations satisfying
$D(x)=[a,x]$, $D(y)=[b,y]$, $D(z)=[c,z]$.
This is the complete infinitesimal envelope of loop preservation and
commutation with $\mathfrak S_3$.  Congruence conditions that do not
descend to this quotient are not silently inserted into its definition.
In particular, $W^{\mathrm{lac}}_n$ is not notation for an unproved
surjectivity statement about the associated graded of the full
profinite solution group.

\begin{claim}{B.g2}{Proved}{The missing special identity}
For $n>1$, normalized representatives identify $W^{\mathrm{lac}}_n$
with the polynomials satisfying \eqref{eq:graded-two},
\eqref{eq:graded-three}, and
\begin{equation}
 J(\psi):=[x,\psi(x,y)]-[x+y,\psi(-x-y,y)]=0.
 \label{eq:graded-special}
\end{equation}
Indeed antisymmetry and \eqref{eq:graded-special} already imply
\eqref{eq:graded-three} in this normalized setting.
\end{claim}
% anchor:B.g2
\manuscriptanchor{\S 42; folios 331, 349; cote 140-3 scans 85, 120}

\begin{proof}
Subtract $\operatorname{ad}a$ from a special derivation to make its
value on $x$ zero.  The remaining value on $y$ has the form
$[y,\psi]$, uniquely in degree $n>1$.
Let $s(x)=y,s(y)=x$ and $r(x)=y,r(y)=z,r(z)=x$.
The difference $sD_\psi s^{-1}-D_\psi$ is inner precisely when
$\psi(y,x)=-\psi(x,y)$: its value on $y$ forces the inner conjugator
to be $\psi$, and its value on $x$ gives this equation.
Similarly
\[
 (rD_\psi r^{-1})(x)=-[z,\psi(y,z)],\qquad
 (rD_\psi r^{-1})(y)=0.
\]
The difference is $\operatorname{ad}\psi$ if and only if
$[x,\psi(x,y)]=[z,\psi(y,z)]$, which by antisymmetry is
\eqref{eq:graded-special}.  Adding its three cyclic transforms gives
$\operatorname{ad}(\psi+r\psi+r^2\psi)=0$.
The free Lie algebra has zero centre, so this is
\eqref{eq:graded-three}.
Conversely this same computation gives invariance under $s$ and $r$.
Moreover
$[y,\psi]=-[z,\psi+r\psi]$, so $D_\psi(z)=-[y,\psi]$
is tangential at $z$.  Thus the derivation is special as well.
\end{proof}

\begin{claim}{B.g3}{Proved}{Exact dimensions through weight twelve}
The independently recomputed dimensions are as follows.  The column
$\mathcal H_n$ is literally the solution space of the two displayed
linear equations, before removing its weight-one line.
\begin{center}
\begin{tabular}{r|rrrr}
$n$&$\dim\mathfrak L_n$&$\dim\mathcal H_n$&$\dim W^{\mathrm{lac}}_n$&$d_n$\\\hline
1&2&1&0&0\\
2&1&0&0&0\\
3&2&1&1&1\\
4&3&1&0&0\\
5&6&2&1&1\\
6&9&3&0&0\\
7&18&6&2&1\\
8&30&10&1&1\\
9&56&19&4&1\\
10&99&33&2&1\\
11&186&62&9&2\\
12&335&112&7&2
\end{tabular}
\end{center}
The review's invariant counts in weights one through eight are correct
for special outer derivations.  Their identification with the solution
space of \eqref{eq:graded-two}--\eqref{eq:graded-three} alone is false.
The first excess over $d_n$ is weight four for the normalized
$\mathcal H_n$, and weight seven for $W^{\mathrm{lac}}_n$.
\end{claim}
% anchor:B.g3
\manuscriptanchor{\S 42; folios 331, 349; cote 140-3 scans 85, 120}

\begin{proof}
Here is a dimension calculation independent of Gaussian elimination.
Let $V$ be the standard two-dimensional representation of
$\mathfrak S_3$ with coordinates $x,y,z$, $x+y+z=0$.
For $g\in\mathfrak S_3$ put
\[
 q_n(g)=\frac1n\sum_{d\mid n}\mu(d)
        \bigl(\operatorname{tr}(g^d\mid V)\bigr)^{n/d}.
\]
The Witt character formula gives
$q_n(g)=\operatorname{tr}(g\mid\operatorname{Lie}_nV)$.
It follows by taking the logarithm of the tensor-algebra character
identity, or equivalently the PBW product, and applying M\"obius
inversion.  The traces of $1,s,r$ on $V$ are $2,0,-1$.
In a standard representation the $s$-odd subspace is one-dimensional
and $1+r+r^2$ vanishes; in the sign and trivial representations the
cyclic sum does not vanish.  Hence
\[
 \dim\mathcal H_n=\frac{q_n(1)-q_n(r)}3.
\]

For $n>1$ the conjugator triples form
$\mathfrak L_n\otimes k\{x,y,z\}$, and the map
$(a,b,c)\mapsto[a,x]+[b,y]+[c,z]$ onto $\mathfrak L_{n+1}$ is
surjective because $[\mathfrak L_n,x]+[\mathfrak L_n,y]
=\mathfrak L_{n+1}$.
There is no conjugator ambiguity in degree greater than one.
Dividing its kernel by the diagonal inner copy of $\mathfrak L_n$
gives, as a virtual $\mathfrak S_3$-representation,
\[
 \operatorname{OutSDer}_n=V\otimes\mathfrak L_n-\mathfrak L_{n+1}.
\]
Taking invariants, exact in characteristic zero, gives
\[
 \dim W^{\mathrm{lac}}_n
 =\frac{q_n(1)-q_n(r)}3
 -\frac{q_{n+1}(1)+3q_{n+1}(s)+2q_{n+1}(r)}6.
\]
Degree one is handled separately: all special derivations there are
inner.  Finally PBW for one generator in every odd weight at least
three gives
$\prod_{n\geq1}(1-t^n)^{-d_n}=(1-t^2)/(1-t^2-t^3)$.
Substitution yields the table.

The program \texttt{graded\_dimensions.py} separately constructs
standard Lyndon brackets in the free associative algebra.  It forms
the exact symmetry projection
$(1-s)(2-r-r^2)$, extracts an independent rational basis, and computes
the rational nullspace of \eqref{eq:graded-special} on it.
A nonzero minor modulo $1000003$ proves independence; the character
formula supplies the matching characteristic-zero upper bound, so
there is no unproved modular-rank inference.  The final kernel
calculation uses exact rational arithmetic, and agrees in every weight
with the formula above.  The code, full output, CSV and integer witness
coefficients are included.
\end{proof}

Write $b_w$ for the standard Lyndon bracket of a word $w$ in $x<y$;
thus $b_x=x,b_y=y$, and $b_w=[b_u,b_v]$ when $v$ is the longest
proper Lyndon suffix and $w=uv$.

\begin{claim}{B.g4}{Refuted}{A fourth-weight obstruction to the naive linearization}
The polynomial
\[
 q_4=b_{xxxy}-4b_{xxyy}+b_{xyyy}
\]
satisfies the two-cycle and three-cycle linear equations but is not a
leading term of any solution of the genuine group equations (I),(II)
over a characteristic-zero field.  The failure already occurs in
weight five.
\end{claim}
% anchor:B.g4
\manuscriptanchor{\S 42; folios 331, 349; cote 140-3 scans 85, 120}

\begin{proof}
Direct substitution verifies both linear equations.  Its special
defect is
\[
 J(q_4)=2b_{xxxxy}-3b_{xxxyy}+b_{xxyxy}
       -3b_{xxyyy}-2b_{xyxyy}+2b_{xyyyy}\ne0.
\]
Thus B.g1--B.g2 already obstruct an actual loop-and-symmetry lift.
For a direct check in Drinfeld's group equations, write
$f=\exp(q_4+q_5+\cdots)$ and
$\log Z=z_0+z_2+\cdots$, where
$z_0=-x-y$ and $z_2=-[x,y]/2$.
Through weight five the equation
$f(Z,X)f(Y,Z)f(X,Y)=1$ is linear in $\log f$ because products of
terms of weight four have weight at least eight.
Its weight-five inhomogeneous term is
\[
 R_5=(\partial_1q_4)(z_0,x)z_2
       +(\partial_2q_4)(y,z_0)z_2.
\]
Consequently $q_5\in\mathfrak L_5$ must satisfy
$q_5+s q_5=0$ and $(1+r+r^2)q_5=-R_5$.
The supplied exact program \texttt{graded\_first\_lifts.py} constructs
this six-unknown system: its coefficient matrix has rank four, whereas
its augmented matrix has rank five.  The output includes a left-kernel
vector annihilating the matrix and pairing to one with the right-hand
side, a complete rational inconsistency certificate.  Thus no higher
coefficients can repair this leading term.
\end{proof}

\begin{claim}{B.g5}{Proved}{An explicit seventh-weight loop solution failing the pentagon}
The polynomial
\begin{align*}
 q_7={}&3b_{xxxyxyy}+3b_{xxxyyxy}+b_{xxyxyxy}
       -3b_{xxyxyyy}\\
      &-3b_{xxyyxyy}-3b_{xxyyyxy}-b_{xyxyxyy}
\end{align*}
belongs to $W^{\mathrm{lac}}_7$ but does not satisfy the infinitesimal
pentagon.  This is an explicit rational graded counterexample; it is
not asserted to be an element of $\GTz$ or of its pro-$\ell$ image.
Moreover $\dim W_7=1$ follows without the 2026 all-weight upper bound.
\end{claim}
% anchor:B.g5
\manuscriptanchor{\S 42; folios 331, 349; cote 140-3 scans 85, 120}

\begin{proof}
The polynomial is an integer linear combination of two independently
computed basis vectors of $W^{\mathrm{lac}}_7$.  The program expands
it in associative words and verifies all three displayed linear
identities exactly.
For a certified nonzero pentagon defect use the following representation
of $\mathfrak t_4$ on $\operatorname{Lie}\langle a,b,c\rangle$.
For $i,j\leq3$, $t_{ij}$ is the derivation sending
$a_i\mapsto[a_i,a_j]$, $a_j\mapsto[a_j,a_i]$ and the third generator
to zero, where $(a_1,a_2,a_3)=(a,b,c)$; $t_{i4}$ acts by
$\operatorname{ad}a_i$.
All defining relations hold: for example the derivation $t_{ij}$ fixes
$a_i+a_j$, so it commutes with
$\operatorname{ad}(a_i+a_j)$; the remaining relations follow by Jacobi
or disjoint support.  The program also checks all eighteen such
relations directly.

Let $P(\psi)$ be the left side minus the right side of the pentagon
(1.4) in the prompt.  Under this representation, the coefficient of
the associative word $a^5bcb$ in $P(q_7)(a)$ is $16$.
It is therefore nonzero in characteristic zero, and hence the original
pentagon defect in $\mathfrak t_4$ is nonzero.  No faithfulness of the
representation is needed for this implication.
\nolinkurl{graded_pentagon_witness.py} prints this coefficient and
writes the entire image; its exact output is part of the archive.
Since $W_7\subset W^{\mathrm{lac}}_7$ by B.g8, a nonzero pentagon
map on the latter two-dimensional space leaves a kernel of dimension
at most one.  The classical free odd-weight subalgebra of B.g7
supplies dimension at least one.  Thus $\dim W_7=1$.
\end{proof}

\subsection{Comparison with the Grothendieck--Teichm\"uller Lie algebra}

\begin{claim}{B.g6}{Known}{The typed rational comparison chain}
Let $\mathrm{GT}^{\mathrm{eq}}_\ell$ mean the group of normalized
invertible pairs satisfying the pro-$\ell$ forms of (I),(II),(III).
This notation is distinct from the image of $\GT$ in that group.
There are continuous homomorphisms
\[
 G_{\mathbb Q}\hookrightarrow\GT
 \longrightarrow\mathrm{GT}^{\mathrm{eq}}_\ell
 \longrightarrow\mathrm{GT}(\mathbb Q_\ell).
\]
The last arrow is injective for normalized free-group pairs; no
surjectivity of either completion arrow is asserted.  On the
characteristic-zero algebraic group with $\lambda=1$, a choice of
associator gives
\[
 \mathrm{GRT}_1(k)\simeq\mathrm{GT}_1(k),\qquad
 \operatorname{gr}\mathfrak{gt}_1(k)=\mathfrak{grt}_1(k).
\]
The first isomorphism depends on the associator, while the induced
map on the graded Lie algebra is the identity in the standard
coordinates.  The space $W$ in the prompt is $\mathfrak{grt}_1$.
\end{claim}
% anchor:B.g6
\manuscriptanchor{\S 42; folios 331, 349; cote 140-3 scans 85, 120}

\begin{proof}
The arithmetic inclusion and the normalized pair convention are
those of \citet[Proposition~1, Proposition~2 and Main Theorem]{HS2000}:
Belyi supplies injectivity of the outer tripod action, and the
normalized representative is unique.  Projection of the free
profinite group to its characteristic pro-$\ell$ quotient sends each
defining relation to the corresponding pro-$\ell$ relation and sends
an automorphism and its inverse to inverse maps.  This defines the
middle arrow; it does not show that every solution in the new category
comes from a profinite one.

The continuous unipotent comparison is
\citet[Theorem~A.6 and Corollary~A.10]{HM2003}.
Its hypotheses hold for the finitely generated discrete free group
$F_2$.  The Magnus expansion embeds its free pro-$\ell$ completion
in $1+\mathbb Z_\ell\langle\!\langle u,v\rangle\!\rangle$.
The expansion can be recovered from all its finite-degree truncations,
and scalar extension to $\mathbb Q_\ell$ loses no coefficient.
Consequently the free pro-$\ell$ group embeds in its unipotent
completion.  Equality of two normalized pairs after completion
therefore gives equality of their scalar parameters and of $f$;
this proves the stated injectivity.

Drinfeld defines the pro-unipotent group in \S4 of
\citet{Drinfeld1990g}.  His Propositions~5.1 and~5.5 give the left
$\mathrm{GT}(k)$-torsor and right $\mathrm{GRT}(k)$-torsor of
associators.  Proposition~5.6 is the exact theorem identifying these
groups after an associator has been chosen; it also explicitly states
$\operatorname{gr}\mathfrak{gt}_1=\mathfrak{grt}_1$ and that the
induced graded map is the identity.  The Lie equations are
(5.17)--(5.20), and Proposition~5.7 removes the redundant special
identity.  Thus these are comparisons of proalgebraic groups, their
Lie algebras and their weight filtrations, not comparisons of arbitrary
finite congruence quotients.
\end{proof}

\begin{claim}{B.g7}{Known}{The arithmetic lower bound and its hypotheses}
Let $I^n_\ell G_{\mathbb Q}$ be the kernel of the outer action on
$F^{(\ell)}/L^{n+1}F^{(\ell)}$.  The positive graded arithmetic Lie
image, tensored with $\mathbb Q_\ell$, contains the free odd-weight
Lie algebra and therefore has dimension at least $d_n$ in weight $n$.
In the standard mixed-Tate comparison this image is itself freely
generated by one element in each odd weight at least three.  Hence
any correctly derived graded condition on the Galois image must leave
at least $d_n$ independent solutions.  Our table satisfies this test
in every computed weight.
\end{claim}
% anchor:B.g7
\manuscriptanchor{\S 42; folios 331, 349; cote 140-3 scans 85, 120}

\begin{proof}
The exact results needed from
\citet{HM2003}, \emph{Weighted completion of Galois groups and Galois
actions on the fundamental group of $\mathbb P^1-\{0,1,\infty\}$},
are Theorem~7.3, Theorem~8.4, Lemma~8.5 and Remark~8.6.
The action factors through the Galois group unramified outside $\ell$,
as used explicitly in their \S8.4.
For this curve $H^1_{\mathrm{et}}(X,\mathbb Q_\ell(1))$ is a trivial
Galois module because the three punctures are rational.  Thus
Theorem~8.4 applies: its graded outer Lie image is precisely the
rationalization of $\operatorname{gr}^{>0}_{I_\ell}G_{\mathbb Q}$.
Theorem~7.3 supplies generators in odd weights, initially including
weight one; Lemma~8.5 kills weight one in the outer action, and
Remark~8.6 identifies the nonzero higher odd abelian classes by
Soul\'e's characters.  Hain--Matsumoto alone proves generation, not
freeness, a limitation stated explicitly in their introduction.

Brown's \emph{Mixed Tate motives over $\mathbb Z$}, Theorem~1.1 and
its stated consequence called Conjecture~1 in the introduction,
prove faithfulness of the mixed Tate action
\citep{Brown2012g}.  The unipotent motivic Lie algebra is free on one
generator in each odd weight at least three, as recalled in that
paper's \S1.1.  Its action respects the braid associativity and
symmetry operations, so its image satisfies the pentagon and the two
cycle equations.  This gives the free odd-weight subalgebra of $W$.
For the $\ell$-adic arithmetic comparison the same generators are
identified by their nonzero Soul\'e abelian classes: the crystalline
mixed Tate completion in Hain--Matsumoto, Theorem~9.2, Proposition~9.6 and
Theorem~9.8, has the same free graded source with weight one removed.
The realization comparison on abelian generators is the regulator
$K_{2n-1}(\mathbb Z)\otimes\mathbb Q_\ell\to H^1_f(G_{\mathbb Q},\mathbb Q_\ell(n))$.
Its isomorphism is the final assertion of Theorem~9.2, with the
crystalline modification stated immediately after Proposition~9.6;
the proof explicitly includes $\ell=2$ via Rognes--Weibel, and odd
primes via Soul\'e.  Thus it is an isomorphism on these one-dimensional
odd abelian generators.  A graded map between free Lie algebras on
these generators that is invertible on their abelianizations is
surjective degree by degree; their dimensions agree by PBW, so it is
an isomorphism.  Brown's faithfulness then rules out relations in the
arithmetic image.  This is also the Deligne--Ihara consequence stated
at the beginning of Brown's paper.  Applying an actual graded
condition to this embedded free subalgebra proves the inequality.
This argument uses positive word weight, not Hain--Matsumoto's
negative doubled motivic weight without reindexing.
\end{proof}

\begin{claim}{B.g8}{Proved}{The special identity follows from the pentagon}
Over any field of characteristic zero, the two-cycle, three-cycle and
infinitesimal pentagon imply \eqref{eq:graded-special}.  Thus this
part of the loop system is of type (a) in the requested classification;
it does not depend on the 2026 upper bound.
\end{claim}
% anchor:B.g8
\manuscriptanchor{\S 42; folios 331, 349; cote 140-3 scans 85, 120}

\begin{proof}
We give the reduction rather than only invoke Drinfeld's
Proposition~5.7.  In degree greater than two forget strand $2$ in
$\mathfrak t_4$ and write
$a=t_{12},b=t_{23},c=t_{24},u=t_{13},v=t_{34}$.
The kernel is the free Lie algebra on $a,b,c$; the base acts by the
usual braid derivations.  This follows directly from the presentation:
$t_{ij}$ acts on the free fibre generators $s_i,s_j$ by
$[s_i,s_j],[s_j,s_i]$ and kills the others; these derivations satisfy
all the base braid relations.
Put $S=a+b+c$ and
\[
 F_a=\psi(u+b,v)-\psi(u,v),\quad
 F_b=\psi(u+a,v+c)-\psi(u,v),\quad
 F_c=\psi(u,v+b)-\psi(u,v).
\]
All three belong to this free kernel.  The braid relations give
\[
 [S,F_b]=[a,\Delta_a]+[c,\Delta_c],\qquad
 \Delta_a=F_b-F_a,\quad\Delta_c=F_b-F_c.
\]
The pentagon and its image under exchange of strands $1,4$ give
\begin{align*}
 \Delta_a&=\psi(a,b+c)-\psi(a,b)-\psi(b,-b-c),\\
 \Delta_c&=-\psi(c,a+b)+\psi(c,b)+\psi(b,-b-a).
\end{align*}
Both vanish at $S=0$.  In that quotient the actions of $u,v$ are
$\operatorname{ad}c,\operatorname{ad}a$; hence
$[b,F_b+\psi(c,a)]=0$.
The degree-$n>1$ centralizer of $b$ is zero, so
$F_b|_{S=0}=-\psi(c,a)$.

Every degree-$n$ Lie polynomial linear in a free direction $T$ and
otherwise involving $a$ is a scalar multiple of
$\operatorname{ad}_a^{n-1}T$.  Write
$\partial_2\psi(a,-a)T=k\operatorname{ad}_a^{n-1}T$.
Differentiating $\psi(a,-a)=0$ gives the same expression for
$\partial_1\psi(a,-a)T$.
Now differentiate the displayed free-kernel identity with
$c=S-a-b$, at $S=0$ in direction $T$.  Substitution of the two
expressions for $\Delta$ gives
\[
 -[T,\psi(c,a)]
 =k(\operatorname{ad}_a^n-\operatorname{ad}_c^n)T
  +[a,\partial_2\psi(b,a)T]-[c,\partial_2\psi(b,c)T].
\]
Set $G(a,b)=[a,\psi(b,a)]+[c,\psi(b,c)]$, $c=-a-b$.
The three-cycle and antisymmetry turn its derivative into
\[
 d_aG(T)=-k(\operatorname{ad}_a^n-\operatorname{ad}_c^n)T.
\]
Equality of mixed derivatives in directions $U,V$, specialized to
$a=a_0,b=-a_0,c=0,U=a_0,V=T$, gives
$(n+1)k\operatorname{ad}_{a_0}^nT=0$.
Indeed the two derivatives of $\operatorname{ad}_a^n$ are respectively
$n\operatorname{ad}_{a_0}^nT$ and
$-\operatorname{ad}_{a_0}^nT$, while the derivatives at $c=0$ vanish.
Thus $k=0$.  All $a$-derivatives of $G$ vanish; characteristic zero and
evaluation at $a=0$ give $G=0$.
Using antisymmetry, this is $J(\psi)=0$.
Weights one and two vanish for the full system: the degree-one
pentagon forces both coefficients to vanish, and the degree-two
three-cycle is $3[x,y]$ times its scalar coefficient.
This completes the proof, including the low weights.  The free-fibre
calculation is also the argument of OpenAI, Lemma~2.4; its logic has
been reproduced here independently of their all-weight estimate.
\end{proof}

\begin{claim}{B.g9}{Proved}{The conditional consequence of the 2026 upper bound}
Assume \citet[Corollary~5.4]{OpenAI2026g}, namely $\dim W_n\leq d_n$ in every
weight.  Then the free arithmetic graded subalgebra in B.g7 equals
$W\otimes\mathbb Q_\ell$ in every weight.  Consequently every
homogeneous linear relation on that Lie algebra which holds on the
arithmetic graded image holds on all of $W\otimes\mathbb Q_\ell$.
When its premise is supplied by the attached manuscript, the equality
is labelled \emph{claimed (OpenAI, 2026, unrefereed)}.
It does not identify $G_{\mathbb Q}$ with $\GT$.
\end{claim}
% anchor:B.g9
\manuscriptanchor{\S 42; folios 331, 349; cote 140-3 scans 85, 120}

\begin{proof}
The injected free subalgebra has dimension $d_n$ in degree $n$ by
PBW.  The opposite inequality forces equality of each finite-dimensional
weight piece.  Taking their products gives equality after weight
completion because each coefficient of a bracket involves finitely
many weight pieces.  A linear relation is a linear map out of such a
piece; a map vanishing on a spanning subspace vanishes on the whole
space.  For a claimed nonlinear condition one would additionally
need a specified algebraic equation and the appropriate Zariski-density
argument; the existence of a tangent space by itself is insufficient.
A filtered or graded Lie algebra does not recover the finite,
disconnected, torsion or integral data lost under completion, so none
of these arguments imply the profinite assertion.
\end{proof}

\paragraph{Audit of the imported upper bound.}
We opened the attached 41-page manuscript and followed the chain from
\S2 to Corollary~5.4, rather than using its abstract as a theorem.
We also opened the attached manuscript of \citet{Willwacher2026g}, its introductory
statements and the beginning of its alternative construction.  The
following distinctions record the extent of our check.

\begin{enumerate}[leftmargin=*]
\item \emph{Integral reduction (OpenAI, Lemmas 2.1--2.4 and
Proposition 2.5).}  The centralizer, free-fibre and derivative argument
for the special identity is reproduced in B.g8.  A rational solution
space intersects the finite free $\mathbb Z_{(2)}$ Lie lattice in a
saturated submodule: if $2v$ belongs, then $v$ belongs to the rational
solution space.  Thus $L\cap2K=2L$, and reduction embeds with dimension
unchanged.  Integral Lyndon bases and the iterated free-fibre PBW basis
make the two equation targets torsion free.  These are exactly the
hypotheses required to reduce the rational equations, and we checked
them.  No claim that the reduction equals \emph{every} solution in
characteristic two is used.
\item \emph{The alphabet $A=x^2,C=[x,y],B=y$ (Lemma 3.1).}
With $x>y$ the leading words $xx,xy,y$ form a prefix code, proving
associative injectivity.  Lazard elimination uses
$e_j=\operatorname{ad}_x^j y$; in characteristic two
$e_{2i}=\operatorname{ad}_A^iB$ and
$e_{2i+1}=\operatorname{ad}_A^iC$.  We checked the inference that every
ordinary Lie element of weight greater than one has this expression.
It does not identify $x^2$ itself as an ordinary Lie element in $x,y$.
\item \emph{The transfer of the pentagon (Lemma 3.3 and
Proposition 3.2).}  We independently evaluated all 63 source
commutator cases needed for the incident and nonincident index patterns
in the free associative fibre over $\mathbb F_2$; all agree with the
displayed formulas.  The ordering argument uses the decreasing pair
(number of base letters, number of base--fibre inversions), so its
termination does not require preservation of symbolic length.
We checked why source corrections increase the $u$-count, target
corrections decrease it, and the common quadratic parts agree.
The independent source normal forms come from the prefix code
$XX,XU_i,U_i$.  These are the logical points on which the claimed
transfer rests.  We have not supplied a second all-weight construction
of the complete target rewriting calculation: the transferred
Proposition~3.2 remains an imported assertion in our conditional use
of Corollary~5.4, \emph{claimed, steps not checked by us} beyond the
checks just specified.
\item \emph{Deletion and injectivity (\S4).}  The integer edge
valuations use the distinct rational hyperplanes $z_i-z_j=0$ and
$z_i+z_j=0$ before reduction.  They must not be merged in characteristic
two.  We checked the logarithmic triangle/residue argument for the
incidence relations, including repeated letters.  Independently,
\nolinkurl{graded_preprint_audit.py} finds the thirteen codimension-two
flats of the $B_3$ arrangement and verifies 13,104 exact integral
deletion incidence relations on all words of lengths zero through
five.  We followed the minimal-$A$ argument: an $e$ before a nonempty
$p,0$ block cannot be deleted before that block, and its final $p$ or
$0$ has coefficient zero modulo two.  Reversal excludes the remaining
terminal-$e$ words.  In Proposition~4.4 the two-marked-position test
forces all $A,B$ coefficients to agree, and the three-marked-position
test kills their common value when there are at least two $A$'s.
We checked separately the one-$A$ exceptions $r=1,2$.
These observations check the stated argument conditional on the
transferred relation; bounded deletion tests alone are not an
all-weight proof.
\item \emph{The image bound (Lemmas 5.1--5.2, Proposition 5.3).}
The $e_j$ expansion makes the number of even elimination letters equal
to the $B$-count.  The Ore algebra with $DB=C,DC=0,X^2=0$ is valid
because $D^2=0$ in characteristic two.  We checked the lower-count
argument proving $f(C,B+s(C))=f(C,B)$.
The subsequent Vandermonde matrix is inverted over
$\mathbb F_2(t_0,\ldots,t_{r-1})$, where its determinant is a nonzero
polynomial; it is not evaluated on the two-element field, and there
is no division by the length $r$.  Divisibility by every
$s_{i-1}-s_i$ excludes $g_0=B$ from the associative expansion in
$g_j=\operatorname{ad}_C^jB$; the retraction $g_0\mapsto0$ preserves
ordinary Lie membership.  We checked this logical passage and the
finite filtration sum in Corollary~5.4.  It remains conditional here
on the preceding imported transfer.
\item \emph{Material not imported as checked.}  OpenAI's holonomy
construction in \S\S6--8, especially Proposition~7.1, is
\emph{claimed, steps not checked by us}; Brown supplies the independent
lower bound used here.  Willwacher's Theorems~1.1 and~1.3 are
\emph{claimed, steps not checked by us}: we have not checked his
complete alternative proof.  The attached Lean comparator ends with
\texttt{theorem main : MainStatement := by sorry}; it is a precise
statement, not an attached proof certificate.  The separate scope
note does not change this fact.
\end{enumerate}

\subsection{Torsion, lifting and what a congruence condition can retain}

\begin{claim}{B.g10}{Proved}{Torsion in the normalized pro-$\ell$ kernel}
The subgroup $\lambda=1$ of normalized loop-and-symmetry automorphisms
of the free pro-$\ell$ group has no nonidentity element of finite
order.  The same is true in characteristic-zero unipotent completion.
For an element of finite order in the full pro-$\ell$ solution group,
its order equals the order of its scalar image in
$\mathbb Z_\ell^\times$.  Thus its possible order divides $\ell-1$
for odd $\ell$, and divides two for $\ell=2$.
Every divisor of $\ell-1$ is realized in the pro-$\ell$ arithmetic
image when $\ell$ is odd.  Finite subgroups of a fixed order in the
full pro-$\ell$ solution group are then conjugate by its $\lambda=1$
kernel.  For $\ell=2$ only the trivial group and order two occur;
we make no conjugacy assertion at two.  The pair $(-1,1)$ has order two.
The odd-prime existence statement does not give explicit normalized
coefficients $f$.
\end{claim}
% anchor:B.g10
\manuscriptanchor{\S 42; folios 331, 349; cote 140-3 scans 85, 120}

\begin{proof}
The normalized representative is unique by the free-group centralizer
argument used in B.g1, so a finite-order outer element has a finite-order
normalized representative.  Suppose $\alpha\ne1$ acts trivially on
abelianization.  The Magnus filtration is separated; take its first
nonzero derivation term $D$ in positive degree.  Composition adds this
term, so the first term of $\alpha^m$ is $mD$.
The module of homogeneous derivations embeds in a finite direct sum
of free $\mathbb Z_\ell$ Lie pieces and is torsion free.  Hence
$mD\ne0$ for every nonzero integer $m$, contradiction.
The same proof works over $\mathbb Q_\ell$, or follows from
exponential and logarithm for a prounipotent group in characteristic
zero.
If an element $g$ has finite order and its scalar has order $e$,
then $g^e$ is a finite-order element in the torsion-free kernel, so
$g^e=1$.  The reverse divisibility is immediate.  The torsion units of
$\mathbb Z_\ell^\times$ are $\mu_{\ell-1}$ for odd $\ell$ and
$\{\pm1\}$ for two.  Finally $X\mapsto X^{-1},Y\mapsto Y^{-1}$
is an involution and represents $(-1,1)$ in the normalized convention.
For completeness, the normalized $\lambda=1$ kernel is a pro-$\ell$
group.  The kernel of the action of $\operatorname{Aut}(F^{(\ell)})$
on the Frattini quotient is pro-$\ell$: on each successive finite
lower exponent-$\ell$ central quotient, automorphisms trivial on the
preceding quotient differ by homomorphisms to an elementary abelian
$\ell$-group.  Taking inverse limits proves the assertion, and the
normalized kernel is a closed subgroup of it.
The cyclotomic character of $G_{\mathbb Q}$ is onto
$\mathbb Z_\ell^\times$.  For odd $\ell$, the inverse image of
$\mu_{\ell-1}$ in the arithmetic image is therefore an extension
of this finite cyclic group by a pro-$\ell$ group.
It has a complement, unique up to conjugation by that kernel.
Here is the requisite form of Schur--Zassenhaus: for a finite
$\ell$-group kernel, induct on its order using a characteristic
elementary abelian subgroup of its centre; on that subgroup the
obstruction and ambiguity are $H^2(C_m,V)$ and $H^1(C_m,V)$, both
zero since multiplication by $m$ is invertible and averaging contracts
the positive-degree cochain complex.  Induction gives existence and
conjugacy.  Passage to the inverse limit follows by compactness of
the finite nonempty sets of complements and conjugators.
The same argument in the full solution group conjugates any two
finite subgroups with the same scalar image.  A complement for
$\mu_{\ell-1}$ contains its subgroups of every divisor order.
At two the displayed complex-conjugation pair supplies the nontrivial
allowed order; no coprime argument applies there.
Finite nilpotent or finite congruence quotients may acquire torsion
because their graded pieces are finite; the torsion-freeness argument
does not apply to those quotients.  Conversely, these separate
pro-$\ell$ results do not classify torsion in the full profinite
group unless a further joint injectivity theorem is supplied.
\end{proof}

\begin{claim}{B.g11}{Open}{The remaining infinite lifting problem, with a verified finite jet}
For the explicit $q_7$ of B.g5, we construct
\[
 F=q_7+F_9+F_{11},\qquad f=\exp F,
\]
satisfying the genuine group equations (I),(II) with $\lambda=1$
through weight twelve.  The complete rational Lyndon coefficients
are supplied: $F_9$ has 46 nonzero coefficients with common denominator
dividing $720$, and $F_{11}$ has 160 with common denominator dividing
$8640$; the even corrections through weight twelve are zero.
Reduction gives a solution in the finite nilpotent Lazard group
\[
 \exp\bigl(\operatorname{Lie}_{\mathbb F_{13}}\langle x,y\rangle/
                 \mathfrak L_{>12}\bigr),
 \qquad |G|=13^{747}.
\]
Its leading seventh-weight infinitesimal pentagon defect remains
nonzero.  We have not settled whether $q_7$ admits compatible
coefficients in an actual free pro-$\ell$ group satisfying (I),(II)
in every degree, or whether any such element lifts to $\GTz$.
The first uncomputed rational correction is weight thirteen; for
$\ell=13$ this is also the first class outside the class-less-than-prime
Lazard argument used for the finite example.
\end{claim}
% anchor:B.g11
\manuscriptanchor{\S 42; folios 331, 349; cote 140-3 scans 85, 120}

\begin{proof}[Proof of the finite assertion and record of the obstruction]
Put $Z=(XY)^{-1}$ and substitute the actual
$\log Z=-\operatorname{BCH}(x,y)$, not $-x-y$.
Since $F$ starts in degree seven, every commutator between two values
of $F$ has degree at least fourteen.  Thus through degree twelve the
logarithm of (II) is exactly
\[
 F(\log Z,x)+F(y,\log Z)+F(x,y).
\]
Equation (I) is likewise $F(x,y)+F(y,x)=0$ in this range.
The program \nolinkurl{graded_rational_jets.py} computes the required
BCH terms, imposes the additional simplifying ansatz that even
corrections vanish, and solves all the remaining exact rational
constraints while retaining the free parameters from earlier weights.
There are 242 unknown Lyndon coefficients, and the full system has
rank 229.  It then substitutes the resulting $F_9,F_{11}$ back into
both identities and checks every associative coefficient through
weight twelve.  Both defects are zero.

The retained failed greedy computation is relevant to the claimed
``successive approximation'' argument in the prompt: choosing all
free ninth-weight coefficients to be zero after the preceding equations
produces an inconsistent tenth-weight system.  This is not an
obstruction to every lift; preserving those parameters and solving
the coupled constraints gives the successful jet just described.
Thus unrestricted choices at one stage do not prove an integral
infinite lifting theorem.

For the finite example, class $12$ is smaller than $13$, so the
BCH and exponential denominators in this class are units in
$\mathbb F_{13}$.  The displayed rational corrections also have
unit denominators at thirteen.  The truncated free Lie algebra
therefore defines its finite BCH group, and its dimension is
$\sum_{n=1}^{12}\dim\mathfrak L_n=747$.
The independent program \texttt{graded\_lazard13.py} repeats the
substitutions directly modulo thirteen and verifies both defects
are zero.  The normalized substitution has identity linear part;
solving successively by degree gives its inverse, so it is an
automorphism of this finite nilpotent group.  The leading pentagon
coefficient computed in B.g5 reduces to $16\equiv3\pmod{13}$.
This establishes the stated finite equation solution and its
infinitesimal pentagon failure without enumerating $13^{747}$ elements.

An infinite pro-$\ell$ lift requires compatible integral choices at
all further nilpotent and $\ell$-power levels.  Rational exponential
alone does not control denominators in unbounded degree.  Moreover
a pro-$\ell$ solution need not lie in the image of the profinite
solution group.  Those are the precise unproved steps; neither is
supplied by our finite computation or by the additive formal
exponential of B.g14.
\end{proof}

\begin{claim}{B.g12}{Proved}{Failure of literal descent of finite congruence maps}
A homomorphism $\Fhat\to H$ to a finite group factors through the
maximal pro-$\ell$ quotient only if its image is an $\ell$-group;
this condition is also sufficient.  No nontrivial finite quotient of
$\Fhat$ extends as a homomorphism from all the rational points of its
characteristic-zero unipotent completion.  Consequently a congruence
condition imposed through a specified finite matrix quotient has no
automatic informative graded form in the system above.
\end{claim}
% anchor:B.g12
\manuscriptanchor{\S 41; folios315--316; cote140-3 scans54--57}

\begin{proof}
Every continuous finite image of a pro-$\ell$ group is an $\ell$-group,
which proves necessity.  Sufficiency is the universal property of the
maximal pro-$\ell$ quotient.  A characteristic-zero unipotent group,
and its inverse limit, has a unique $m$th root for every positive
integer $m$, given by $\exp((\log g)/m)$.
If a homomorphism to a finite group of exponent $m$ were nontrivial,
write $g=h^m$; its image would nevertheless be one.  Thus every such
homomorphism is trivial.  A finite congruence predicate can still imply
some rational Lie relation, but that requires an explicit comparison
argument.  It cannot be inferred by substituting $z=-x-y$ in the
original finite quotient.
\end{proof}

We found no map identifying the degenerate pentagon (3.5), or a
$2$-adic version of it, with one of Grothendieck's level-two or
level-four congruence conditions.

\begin{claim}{B.g13}{Proved}{The characteristic-two arrangement is not naive reduction of its hyperplanes}
The Lie algebra used in the degenerate pentagon of OpenAI,
Proposition~3.2, retains nine degree-one generators when $d=3$.
The arrangement obtained by identifying the equations
$z_i-z_j=0$ and $z_i+z_j=0$ after reduction to $\mathbb F_2$ has
only six hyperplanes.  Their holonomy Lie algebras are therefore
not identified by the reduction used in that proposition.
\end{claim}
% anchor:B.g13
\manuscriptanchor{\S 38; fol.297; cote140-3 scan18; added Recent results task4}

\begin{proof}
The defining incidence relations are brackets, so impose no
relations in degree one.  Over $\mathbb Q$ the three coordinate
hyperplanes and six signed diagonal hyperplanes are distinct.
Reduction of the integral incidence \emph{presentation} preserves
these nine generator symbols.  Reducing the hyperplane \emph{equations}
first identifies each signed pair, leaving three coordinate and three
diagonal hyperplanes.  The degree-one dimensions, nine and six, are
already different.  Our deletion program accordingly constructs all
thirteen rational codimension-two flats before any scalar reduction.
This checks a concrete obstruction to the proposed naive comparison;
it neither proves nor presupposes a link to a modular congruence
quotient of level two or four.
\end{proof}

\begin{claim}{B.g14}{Proved}{An infinite element in the additive graded formal loop group}
Let $\mathcal A^{\mathrm{add}}$ be the group of continuous outer
$k$-Lie automorphisms of $\widehat{\mathfrak L}$, equal to the identity
on abelianization, preserving the three adjoint-conjugacy classes of
$x,y,z=-x-y$, and commuting with the linear $\mathfrak S_3$ action.
The exponential of $D_{q_7}$ defines a nonidentity element of
$\mathcal A^{\mathrm{add}}$ outside $\mathrm{GRT}_1(k)$.
This supplies an explicit infinite element in this additive graded
formal group.  We do not identify $\mathcal A^{\mathrm{add}}$ with
$\mathrm{GT}_0(k)$ or with any pro-$\ell$ group.
\end{claim}
% anchor:B.g14
\manuscriptanchor{\S 42; folios 331, 349; cote 140-3 scans 85, 120}

\begin{proof}
The derivation raises degree by seven, so
$\exp D_{q_7}=\sum_{j\geq0}D_{q_7}^j/j!$ converges for the weight
filtration, with only finitely many terms affecting a fixed weight.
It is an automorphism, with inverse $\exp(-D_{q_7})$.
Tangentiality integrates directly.  If $D(u)=[u,c]$, set
$U(t)=\exp(tD)u$ and $C(t)=\exp(tD)c$.  Then
$U'(t)=[U(t),C(t)]$.
The equation $g'(t)=g(t)C(t)$, $g(0)=1$, has a unique solution
weight by weight in the completed enveloping algebra, using integration
of polynomials in $t$ over the characteristic-zero field.
The solution is group-like because $C(t)$ is primitive, and
$U(t)=g(t)^{-1}ug(t)$.  Apply this to $x,y,z$.
In the quotient by inner derivations, B.g2 says that $D_{q_7}$ is
$\mathfrak S_3$-invariant.  Exponential commutes with conjugation and
with this quotient, so its outer class centralizes $\mathfrak S_3$.
Finally the degree-seven logarithmic term of an element of
$\mathrm{GRT}_1(k)$ must satisfy the infinitesimal pentagon.
B.g5 proves that $q_7$ does not; higher terms cannot cancel that
lowest defect.  Thus the constructed element is outside
$\mathrm{GRT}_1(k)$, as asserted.
\end{proof}

\begin{claim}{B.g15}{Proved}{A genuine finite non-arithmetic outer solution}
Let $N_2$ be the characteristic finite quotient of $\Fhat$ given by
the free class-twelve BCH group over $\mathbb F_{13}$ in B.g11.
The normalized automorphism defined there solves (I),(II) but is not
the projection of an element of $\GT$.  Its outer class is not in
the image of $G_{\mathbb Q}$ in $\operatorname{Out}(N_2)$.
This gives a finite non-arithmetic equation solution, of order-quotient
$|N_2|=13^{747}$; it does not prove that the solution lifts to $\GTz$.
It is an additional nilpotent quotient, not one of the small matrix
or permutation quotients prescribed in B(e).
\end{claim}
% anchor:B.g15
\manuscriptanchor{\S 42; folio349; cote140-3 scan120}

\begin{proof}
We first supply an actual finite pure-braid target for the pentagon;
replacing it by the exponential of a braid Lie algebra without a
comparison would not suffice.
Forgetting the fourth strand gives
$1\to F_3\to P_4\to P_3\to1$, with fibre generators
$a_i=A_{i4}$, $1\leq i\leq3$.
Conjugation gives an action of $P_4$ on its normal free fibre.
Let
\[
 N_3=\exp\bigl(\operatorname{Lie}_{\mathbb F_{13}}
                 \langle a,b,c\rangle/\mathfrak L_{>8}\bigr).
\]
The fibre action descends to $N_3$.  The full twist is central in
$P_4$, hence acts trivially on the fibre; thus the action factors
through $P_4/Z(P_4)=\pi_1(M_{0,5})$.
This identifies a finite quotient of the actual pentagon target.
The split free-fibre sequence follows directly from forgetting a
point in the ordered configuration space: locally the punctures
move by disjoint-disc isotopies, its fibre is a three-punctured
plane, and induction gives asphericity and the stated fundamental
sequence.  A far-away fourth point supplies a section.

Here are explicit signs for its action.  In the inverse Artin
convention the base half-twists act by
\[
 \beta_i(a_i)=a_{i+1},\qquad
 \beta_i(a_{i+1})=a_{i+1}^{-1}a_i a_{i+1},
\]
fixing the other generator.  The base pure generators are
$\beta_1^2,\ \beta_2\beta_1^2\beta_2^{-1},\ \beta_2^2$.
Their first Lie terms are precisely the derivations $t_{12},t_{13},t_{23}$
used in B.g5.  The fibre generators act by inner conjugation and
have first terms $\operatorname{ad}a_i=t_{i4}$.
The program \texttt{graded\_braid\_transfer.py} verifies the Artin
relation as reduced free words, all six first Lie terms, and the
identity that the base full twist acts by conjugation with
$(abc)^{-1}$; fibre conjugation by $abc$ cancels it.
Thus the signs and central quotient in this transfer are checked
without an associator or a braid-formality assertion.

The image lies in the group of automorphisms inducing the identity
on $H_1(N_3,\mathbb F_{13})$.  If $\operatorname{IA}_i$ denotes
those raising Lie degree by at least $i$, then
$[\operatorname{IA}_i,\operatorname{IA}_j]\subset
\operatorname{IA}_{i+j}$ and $\operatorname{IA}_8=1$.
This finite target has class at most seven.  As a linear operator
$\phi=1+T$ on the truncated Lie algebra, $T^8=0$, whence
$\phi^{13}=1+T^{13}=1$ in characteristic thirteen.
Truncated exponential and logarithm identify these degree-raising
automorphisms with degree-raising derivations.  To check the latter
statement directly, $\exp(t\log\phi)=\phi^t$ for all
$t\in\mathbb F_{13}$ preserves brackets.  In every relevant
homogeneous component the bracket-preservation polynomial has degree
less than thirteen in $t$, so it vanishes identically; differentiating
at zero shows that $\log\phi$ is a derivation.
All denominators used have index at most eight and are units.
Thus the free BCH group $N_2$ admits evaluation homomorphisms into
this target for arbitrary pairs of its elements: apply the universal
property of the free truncated Lie algebra to their logarithms.
The same property shows that the kernel defining $N_2$ is fully
invariant in $\Fhat$, hence characteristic.

Evaluate the genuine group pentagon, equation (4.10) of
\citet{Drinfeld1990g} in this finite pure-braid target.
Every occurrence of $f$ starts in degree seven.  Since the target
has class seven, its value depends only on $q_7$ and on the
degree-one parts of the braid arguments.  Products of braid
arguments give the sums of chords in (1.4), so the lowest group
pentagon defect is exactly the infinitesimal defect evaluated in
B.g5.  On $a$ its coefficient at the word $a^5bcb$ is
$16\equiv3\pmod{13}$.  This word has length eight and survives
in $N_3$; the defect is nonidentity.  The corrections $F_9,F_{11}$
vanish in this target and cannot cancel it.  Therefore our finite
pair is not the projection of a pair satisfying the group pentagon.

Finally we check that taking outer classes cannot remove the
obstruction.  In $N_2$,
\[
 C_{N_2}(X)=\langle X\rangle Z(N_2),\quad
 C_{N_2}(Y)=\langle Y\rangle Z(N_2),\quad
 Z(N_2)=\exp\mathfrak L_{12}.
\]
Indeed the associative centralizer of a free letter consists of its
power series, and its intersection with the ordinary free Lie algebra
below the top degree is just the line spanned by that letter.
The invertible factor in $\exp(\operatorname{ad}x)-1
=\operatorname{ad}x\,(1+\cdots)$ identifies the group and Lie
centralizer conditions.  Hence the displayed formulas hold in this
class-less-than-prime quotient.
If two normalized pairs $f,g\in[N_2,N_2]$ give inner-conjugate
automorphisms fixing $X$, their inner conjugator is $X^a z$ with
$z$ central.  Comparing their values on $Y$ gives a difference by
an element of $\langle Y\rangle Z(N_2)$; abelianization forces both
powers of $X,Y$ to be zero.  Thus $f$ and $g$ differ only by
$Z(N_2)$.  That difference is invisible in the class-seven pentagon
target, so both have the same nonzero defect.
Any arithmetic outer class equal to ours has scalar $1$ modulo
thirteen by its action on abelianization, and has such a normalized
representative.  The Galois image lies in $\GT$, so its pentagon
defect must vanish.  This contradiction proves the outer-class
assertion as well.
\end{proof}
